\section{习题与解答}
\begin{exercise}\label{ex:o1oval}
要用一个折射界面把介质$n_1$中的点$A$严格成像到介质$n_2$中的点$B$。忽略反射，证明等光程曲面$n_1|Q-A|+n_2|B-Q|=C$的适当光滑部分满足折射定律，说明与球面近轴成像的区别。
\Diagram{fig:o1ovalq}{笛卡尔卵形面的一个分支，图取$n_1=1,n_2=1.5,C=7.5$}{\fill(0,0)circle(1.5pt)node[below]{$A$};\fill(5.5,0)circle(1.5pt)node[below]{$B$};\draw[curve]plot coordinates {(3.732186,-1.500000) (3.532345,-1.462025) (3.375574,-1.424051) (3.242276,-1.386076) (3.124437,-1.348101) (3.017837,-1.310127) (2.919932,-1.272152) (2.829045,-1.234177) (2.744007,-1.196203) (2.663963,-1.158228) (2.588269,-1.120253) (2.516428,-1.082278) (2.448047,-1.044304) (2.382811,-1.006329) (2.320466,-0.968354) (2.260808,-0.930380) (2.203667,-0.892405) (2.148908,-0.854430) (2.096420,-0.816456) (2.046116,-0.778481) (1.997930,-0.740506) (1.951812,-0.702532) (1.907731,-0.664557) (1.865669,-0.626582) (1.825624,-0.588608) (1.787608,-0.550633) (1.751645,-0.512658) (1.717772,-0.474684) (1.686040,-0.436709) (1.656510,-0.398734) (1.629252,-0.360759) (1.604347,-0.322785) (1.581883,-0.284810) (1.561951,-0.246835) (1.544645,-0.208861) (1.530055,-0.170886) (1.518266,-0.132911) (1.509353,-0.094937) (1.503375,-0.056962) (1.500375,-0.018987) (1.500375,0.018987) (1.503375,0.056962) (1.509353,0.094937) (1.518266,0.132911) (1.530055,0.170886) (1.544645,0.208861) (1.561951,0.246835) (1.581883,0.284810) (1.604347,0.322785) (1.629252,0.360759) (1.656510,0.398734) (1.686040,0.436709) (1.717772,0.474684) (1.751645,0.512658) (1.787608,0.550633) (1.825624,0.588608) (1.865669,0.626582) (1.907731,0.664557) (1.951812,0.702532) (1.997930,0.740506) (2.046116,0.778481) (2.096420,0.816456) (2.148908,0.854430) (2.203667,0.892405) (2.260808,0.930380) (2.320466,0.968354) (2.382811,1.006329) (2.448047,1.044304) (2.516428,1.082278) (2.588269,1.120253) (2.663963,1.158228) (2.744007,1.196203) (2.829045,1.234177) (2.919932,1.272152) (3.017837,1.310127) (3.124437,1.348101) (3.242276,1.386076) (3.375574,1.424051) (3.532345,1.462025) (3.732186,1.500000)};\draw[curve](0,0)--(2.372225,1)--(5.5,0);\draw[guide](0,0)--(2.372225,-1)--(5.5,0);\node[above]at(2.372225,1){$Q$};}
\end{exercise}
\begin{solution}
定义$F(Q)=n_1r_1+n_2r_2$，沿入射传播方向$\vct t_1=(Q-A)/r_1$、出射方向$\vct t_2=(B-Q)/r_2$。梯度为$\nabla F=n_1\vct t_1-n_2\vct t_2$。图\ref{fig:o1ovalq}的等值面法线沿$\nabla F$，任何切向位移$\delta\vct q$满足$\nabla F\cdot\delta\vct q=0$，所以$n_1\vct t_{1\parallel}=n_2\vct t_{2\parallel}$。
\Diagram{fig:o1ovala}{等光程面的法线由入射、出射约化方向矢量之差确定}{\draw[axis](0,0)--(2.4,1.2)node[above left]{$n_1\vct t_1$};\draw[axis](0,0)--(2.1,-.5)node[below]{$n_2\vct t_2$};\draw[->,thick](2.1,-.5)--(2.4,1.2);\node[right]at(2.6,.35){差矢量};\draw[guide](0,0)--(.3,1.7);}
上式为斯涅尔切向条件，法线的矢量构造见图\ref{fig:o1ovala}。须选择使光由介质1进入介质2、无全反射且不被其他界面截断的分支。球面一般只能在轴上小孔径处近似满足固定物像等光程；这个等光程面通常不是球面，可以为指定物像对校正球差，却不保证别的物点也严格成像。
\end{solution}

\begin{exercise}\label{ex:o1galilean}
空气中正物镜焦距$f_1>0$，负目镜$f_2<0$，间隔$d=f_1+f_2>0$。求无焦系统角放大率、轴向平行束宽度比和像的正倒，与双正透镜系统比较。
\Diagram{fig:o1galileoq}{负目镜截取尚未到物镜焦点的会聚光束}{\OptAxis\draw[<->](1,-1.3)--(1,1.3);\draw[>-<](4,-1)--(4,1);\draw[guide](5.5,-1.5)--(5.5,1.5);\node[above]at(5.5,1.5){物镜焦点};\node[above]at(4,1){负目镜};}
\end{exercise}
\begin{solution}
图\ref{fig:o1galileoq}仍用两透镜矩阵，把$d=f_1+f_2$代入得到
\[
 M=\begin{pmatrix}-f_2/f_1&f_1+f_2\\0&-f_1/f_2\end{pmatrix}.
\]
因$f_2<0$，两个对角元都正。轴向平行光的宽度比$|f_2|/f_1$，小视角放大率$\theta_o/\theta_i=f_1/|f_2|>0$，所以角方向不反转。
\Diagram{fig:o1galileoa}{光束在到达几何焦点之前被负透镜重新准直}{\OptAxis\draw[<->](1,-1.3)--(1,1.3);\draw[>-<](4,-1)--(4,1);\draw[->](0,1)--(1,1)--(4,.333333)--(6,.333333);\draw[->](0,-1)--(1,-1)--(4,-.333333)--(6,-.333333);\draw[guide](4,.333333)--(5.5,0);\draw[guide](4,-.333333)--(5.5,0);}
图\ref{fig:o1galileoa}按$f_1=4.5,f_2=-1.5,d=3$作图，束宽缩为1/3、角放大率3，内部没有经过实焦点。双正系统角放大率为负，会产生倒向角像。两系统都满足$\Phi=0$，但角放大率符号和内部光路不同，对应正立、倒立两种观察效果。
\end{solution}

\begin{exercise}\label{ex:o1layers}
平行分层透明介质的各层厚度为$h_j$、折射率$n_j$。底部点物体从顶面外的介质$n_0$中近轴观察。求表观深度，说明为何各层的真实光程厚度不能直接作为表观深度相加。
\Diagram{fig:o1layersq}{平行分层结构保持每条光线的切向约化分量}{\foreach\y in{0,.7,1.4,2.1}{\draw[curve](-2,\y)--(2,\y);}\fill(0,0)circle(1.5pt)node[below]{物点};\draw[->](0,0)--(.35,.7)--(.6,1.4)--(1.1,2.1)--(1.8,3);\node[right]at(2,.35){$n_3,h_3$};\node[right]at(2,1.05){$n_2,h_2$};\node[right]at(2,1.75){$n_1,h_1$};}
\end{exercise}
\begin{solution}
图\ref{fig:o1layersq}所有平面界面无光焦度，$p=n\theta$在近轴下保持。光从底部穿过第$j$层的横向位移为$h_j\theta_j=h_jp/n_j$，到顶面总横移$y=p\sum_jh_j/n_j$。外部出射斜率为$\theta_0=p/n_0$。
\Diagram{fig:o1layersa}{外部光线反向延长至虚像，深度由横移除外部斜率决定}{\draw[curve](-2,1.6)--(2,1.6);\fill(0,0)circle(1.5pt)node[below]{真实物};\fill(0,.6)circle(1.5pt)node[left]{虚像};\draw[guide](0,.6)--(1,1.6);\draw[->](1,1.6)--(2.3,2.9);\draw[guide](0,0)--(0,2.6);}
图\ref{fig:o1layersa}表观深度$h_a=y/\theta_0=n_0\sum_jh_j/n_j$。真实传播光程为$\sum_jn_jh_j$，一个含$n_j$倒数、另一个含$n_j$本身，物理意义不同。若各层折射率相同$n$，便恢复$n_0h/n$；若$n_0=n$，深度不变。
\end{solution}

\begin{exercise}\label{ex:o1immersion}
同一薄双凸透镜折射率$n_l=1.50$，在空气中焦距$f_a$。完全浸入折射率$n_0=1.33$的液体，求主面焦距与$f_a$之比，讨论液体折射率超过玻璃时的光焦度符号。
\Diagram{fig:o1immersionq}{同一曲率透镜置于不同外部介质}{\OptAxis\draw[curve](2,-1.3)..controls(1.4,-.4)and(1.4,.4)..(2,1.3);\draw[curve](2,-1.3)..controls(2.6,-.4)and(2.6,.4)..(2,1.3);\node[above]at(2,1.5){$n_l$};\node at(4,1){$n_0$};}
\end{exercise}
\begin{solution}
令曲率组合$K=1/R_1-1/R_2>0$。图\ref{fig:o1immersionq}在外部介质中约化光焦度$\Phi=(n_l-n_0)K$，两侧相同折射率的主面焦距$f_0=n_0/\Phi$。空气中$f_a=1/[(n_l-1)K]$，所以
\[
 \frac{f_0}{f_a}=\frac{n_0(n_l-1)}{n_l-n_0}
 =\frac{1.33\times.50}{.17}\simeq3.912.
\]
\Diagram{fig:o1immersiona}{外介质接近透镜折射率，折射偏折变弱、焦距增大}{\OptAxis\draw[<->](1.5,-1.2)--(1.5,1.2);\draw[->](0,.8)--(1.5,.8)--(5.8,0);\draw[guide](1.5,.8)--(2.6,0);\node[below]at(2.6,0){空气焦点};\node[below]at(5.8,0){液体焦点};}
图\ref{fig:o1immersiona}焦距增大是相对折射率差减小的结果。$n_0=n_l$时界面不可光学分辨，薄透镜光焦度为零；$n_0>n_l$时$\Phi<0$，相同双凸形状转为发散。透镜的会聚、发散性质由相对折射率和几何形状共同决定。
\end{solution}

\begin{exercise}\label{ex:o1mirage}
平移不变介质中$n(y)=n_A-gy$，$g>0$，只取$n>0$区域。光从$(z,y)=(0,0)$向上方以相对$+z$角$\alpha_0\in(0,\pi/2)$进入。求精确轨迹、最高点和回到$y=0$的水平距离。
\Diagram{fig:o1mirageq}{折射率向上减小，光线弯回高折射率区}{\OptAxis\draw[axis](0,0)--(0,3)node[above]{$y$};\draw[axis](0,0)--(1.4,1)node[above]{$\alpha_0$};\draw[guide](0,2.5)--(5.5,2.5);\node[right]at(5.5,2.5){$n$较小};}
\end{exercise}
\begin{solution}
图\ref{fig:o1mirageq}的守恒量$b=n\cos\alpha=n_A\cos\alpha_0$。轨迹斜率$y'=\tan\alpha$，所以$n^2=b^2(1+y'^2)$。在上升支$y'>0$，令$n=b\cosh u$，则$y'=\sinh u$；由$n=n_A-gy$，$\dd y=-(b/g)\sinh u\dd u$，于是$\dd z=-(b/g)\dd u$。
\Diagram{fig:o1miragea}{轨迹是倒置悬链线的一段，对最高点左右对称}{\OptAxis\draw[axis](0,0)--(0,2.7)node[above]{$y$};\draw[curve,domain=0:2.8,samples=80]plot(\x,{1.4*(cosh(1)-cosh((\x-1.4)/1.4))});\draw[guide](1.4,0)--(1.4,.760313);\node[above]at(1.4,.85){最高点};\node[below]at(2.8,0){返回};}
以最高点$u=0,z=z_t$为参考点，上升支满足
\[
 \int_{z_t}^{z}\dd Z=-\frac bg\int_0^u\dd q,
 \qquad z-z_t=-\frac bg u.
\]
故$u=g(z_t-z)/b$。由初始条件$n(0)=n_A=b\cosh(gz_t/b)$求出$z_t$，再以$\cosh$的偶性接续下降支，得到
\[
 y(z)=\frac{n_A-b\cosh[g(z-z_t)/b]}g,\quad
 z_t=\frac bg\operatorname{arcosh}\frac{n_A}b.
\]
最高点$y_t=(n_A-b)/g$，回到$y=0$时$z=2z_t$。图\ref{fig:o1miragea}按一组比例画出同一悬链线结构。这个解适用于有限角度。前面的$n^2$抛物分布也有精确的正弦型轨迹，但其轴向周期依赖初始方向；仅在近轴极限下才化为固定节距$2\pi a$。两种分布的轨迹都须限制在实际介质覆盖的区域内。
\end{solution}

\begin{exercise}\label{ex:o1chromatic}
薄透镜几何曲率固定，空气中折射率随波长变化。已知$n(\lambda_0)=n_0$，一小段波长变化对应$\Delta n$。求焦距一阶变化及平行白光的轴向色差方向。
\Diagram{fig:o1chromq}{同一透镜对不同波长的光焦度不同}{\OptAxis\draw[<->](1.8,-1.3)--(1.8,1.3);\draw[->](0,1)--(1.8,1);\draw[curve](1.8,1)--(4,0);\draw[guide](1.8,1)--(5.6,0);\node[below]at(4,0){短波};\node[below]at(5.6,0){长波};}
\end{exercise}
\begin{solution}
图\ref{fig:o1chromq}的$f=K^{-1}(n-1)^{-1}$。固定$K$求导，$\dd f/\dd n=-K^{-1}(n-1)^{-2}=-f/(n-1)$，所以$\Delta f\simeq-f_0\Delta n/(n_0-1)$。对题图中的正透镜，正常色散使短波折射率较大，蓝光因而比红光更早会聚；发散透镜也服从同一有向焦距微分式。
\Diagram{fig:o1chroma}{像面固定时，不同焦点在屏上形成不同半径模糊斑}{\OptAxis\draw[<->](1,-1.2)--(1,1.2);\draw[guide](4,-1.6)--(4,1.6);\draw[curve](1,.9)--(3,0)--(5,-.9);\draw[curve](1,-.9)--(3,0)--(5,.9);\draw[guide](1,.9)--(4.5,0);\draw[guide](1,-.9)--(4.5,0);}
图\ref{fig:o1chroma}说明固定屏幕不可能同时处在所有波长的几何焦面。色差是$n(\lambda)$改变光焦度的几何效应，单色衍射则即使焦距相同仍形成有限光斑；两种宽度不能不分机制地直接相加。本一阶式还要求$|\Delta n|\ll n_0-1$。
\end{solution}

\begin{exercise}\label{ex:o1prismminimum}
折射率$n$、顶角$A$的棱镜在空气中，光由一个斜面入射、另一个出射。推导对称路径下的最小偏向角公式，说明该状态需满足的传播条件。
\Diagram{fig:o1prismq}{一般穿透光路；法线与界面垂直，入射、折射角分别对法线量取}{
\draw[curve](0,0)--(2,3.4641)--(4,0)--cycle;
\draw[->](-.7,.395855)--(.692820,1.2);
\draw[curve](.692820,1.2)--(3.175128,1.428721);
\draw[->](3.175128,1.428721)--(4.7,1.190565);
\draw[guide](0,1.6)--(1.385641,.8);
\draw[guide](2.482307,1.028721)--(3.867948,1.828721);
\node at(2,2.75){$A$};\node at(2,.55){$n>1$};
\node[above left]at(.28,1.23){$i_1$};\node[below right]at(.98,1.2){$r_1$};
}
\end{exercise}
\begin{solution}
图\ref{fig:o1prismq}中折射角满足$r_1+r_2=A$，空气侧入、出射角满足$\sin i_j=n\sin r_j$，总偏向为$D=i_1+i_2-A$。令$f(r)=\arcsin(n\sin r)$，则$D=f(r_1)+f(A-r_1)-A$。在$n>1$、$0\leq r_j<\arcsin(1/n)$的通常穿透分支，
\begin{align*}
 f'(r)&=\frac{n\cos r}{\sqrt{1-n^2\sin^2r}},\\
 f''(r)&=-\frac{n\sin r}{\sqrt{1-n^2\sin^2r}}
 +\frac{n^3\sin r\cos^2r}{(1-n^2\sin^2r)^{3/2}}\\
 &=\frac{n(n^2-1)\sin r}{(1-n^2\sin^2r)^{3/2}}\geq0.
\end{align*}
凸性保证$r_1=r_2=A/2$为极小点，光路可逆性则给出对称驻点。
\Diagram{fig:o1prisma}{最小偏向的对称光路；图取顶角60度、折射率1.5}{
\draw[curve](0,0)--(2,3.4641)--(4,0)--cycle;
\draw[->](-.7,.8929)--(.8083,1.4);
\draw[->](.8083,1.4)--(3.1917,1.4);
\draw[->](3.1917,1.4)--(4.7,.8929);
\draw[guide](2,-.2)--(2,3.7);
\draw[guide](.1155,1.8)--(1.5011,1);
\draw[guide](2.4989,1)--(3.8845,1.8);
\node at(1.3,1.1){$r_1$};\node at(2.7,1.1){$r_2$};
\node[below]at(2,-.3){$r_1=r_2=A/2$};
\node at(2,2.85){$A$};
}
图\ref{fig:o1prisma}的对称路径有$i_1=i_2=i$，因此
\[
 D_{\min}=2i-A,
 \qquad n=\frac{\sin[(A+D_{\min})/2]}{\sin(A/2)}.
\]
可传播条件为$n\sin(A/2)\leq1$；等号对应掠入、掠出极限，有限通光通常要求严格小于1。若不满足，无法让光沿这个两面穿透分支通过。薄棱镜$A\ll1$且$nA\ll1$时$i\simeq nA/2$，于是$D_{\min}\simeq(n-1)A$。近似要求两个空气侧角度也小，不能只检查棱镜顶角而忽略很大的折射率。
\end{solution}

\begin{exercise}\label{ex:o1obliqueplate}
平面波以角$i$穿过厚$h$、折射率$n$的平行板，板外为空气。比较板后同一平面上的相位与未插板情况，推导附加相位，避免把折线实际路长误当程差。
\Diagram{fig:o1obliqueplateq}{比较同一输出平面上的波相位}{\draw[curve](-2,0)--(2,0);\draw[curve](-2,1)--(2,1);\draw[->](-1.5,-1)--(-.5,0)--(.1,1)--(1.1,2);\draw[guide](-.5,0)--(.5,1);\node[right]at(2,.5){$h$};}
\end{exercise}
\begin{solution}
图\ref{fig:o1obliqueplateq}中平行界面保持切向波矢，$\sin i=n\sin r$。在固定横坐标比较上下参考面的相位增长，板内法向波矢$k_{zg}=k_0n\cos r$，原空气中为$k_{za}=k_0\cos i$。所以附加相位$\Delta\phi=k_0h(n\cos r-\cos i)$。
\Diagram{fig:o1obliqueplatea}{沿不同几何射线积累的光程须先归到同一参考波面}{\draw[curve](0,0)--(4,0);\draw[curve](0,1)--(4,1);\draw[guide](1,-.5)--(1,1.7);\draw[->](.3,-.6)--(1,0)--(1.5,1)--(2.2,1.6);\node[above]at(1,1.8){固定横坐标};}
图\ref{fig:o1obliqueplatea}的计算自动扣除了横向位移相关相位。正常入射恢复$k_0(n-1)h$。若题目改为固定两个点之间的封闭干涉臂，需同时满足光线连接、元件位置和合束条件，不能只把本平面波相位式套给任意折线路径。
\end{solution}

\begin{exercise}\label{ex:o1matrixinverse}
空气近轴系统矩阵为$M$且$\det M=1$。写出逆矩阵，解释“数学逆传播”与把整套透镜实际前后翻转为何须先统一坐标基底。
\Diagram{fig:o1inverseq}{同一对参考面之间的前向、反向关系}{\draw[curve](1,-1)rectangle(3.5,1);\node at(2.25,0){$M$};\draw[->](-.5,.5)--(.8,.5);\draw[->](3.7,-.5)--(5,-.5);}
\end{exercise}
\begin{solution}
图\ref{fig:o1inverseq}用同一固定$z$和有向倾角描述时，数学逆为$M^{-1}=\bigl(\begin{smallmatrix}D&-B\\-C&A\end{smallmatrix}\bigr)$，因为与原矩阵相乘给单位阵。逆关系把末端的$y,n\theta$恢复到起点，并不自动改变倾角的参考轴。
\Diagram{fig:o1inversea}{反向实际传播若重新令光轴向前，倾角符号也改变}{\draw[axis](0,.7)--(4,.7)node[right]{$+z$};\draw[axis](4,-.7)--(0,-.7)node[left]{$+\tilde z$};\node at(2,0){同一光线需换基};}
若反向传播时重新把轴指向左方而高度仍向上，约化倾角需变号，换基矩阵$K=\operatorname{diag}(1,-1)$，反向系统矩阵为$KM^{-1}K=\bigl(\begin{smallmatrix}D&B\\C&A\end{smallmatrix}\bigr)$。图\ref{fig:o1inversea}说明相同光学系统反向后的$B,C$正负取决于坐标约定，不能只看一处符号判定互易性失效。
\end{solution}

\begin{exercise}\label{ex:o1focusplate}
物点到固定屏距离$D$，中间放焦距$f$的薄透镜，空气中近轴成像。求允许的透镜位置数及放大率乘积，讨论$D<4f$。
\Diagram{fig:o1focusq}{固定物屏间距，透镜位置待定}{\OptAxis\draw[->](0,0)--(0,1);\draw[<->](2,-1.2)--(2,1.2);\draw[curve](6,-1.5)--(6,1.5);\draw[<->](0,-1.8)--(6,-1.8)node[midway,below]{$D$};}
\end{exercise}
\begin{solution}
图\ref{fig:o1focusq}中取物体在0、屏在$D$，透镜位置为$u$，于是像距$v=D-u$。由$1/u+1/v=1/f$，相乘得到$uv=f(u+v)=fD$，即
\[
 u^2-Du+fD=0,\qquad
 u_\pm=\frac{D\pm\sqrt{D^2-4fD}}2.
\]
题中的实物、实屏位于透镜两侧，故要求$0<u<D$；当$D>4f$、$f>0$时两个根均满足此条件。
\Diagram{fig:o1focusa}{两次成像透镜位置关于物屏中点对称}{\OptAxis\draw[<->](1.7,-1.2)--(1.7,1.2);\draw[guide](4.3,-1.2)--(4.3,1.2);\draw[curve](6,-1.5)--(6,1.5);\draw[guide](3,-1.7)--(3,1.7);\node[above]at(1.7,1.2){位置一};\node[above]at(4.3,1.2){位置二};}
图\ref{fig:o1focusa}的两个位置将$u,v$互换，放大率分别为$-v/u$和$-u/v$，乘积1。一幅放大像与一幅缩小像配对出现。两位置间距
\[
 l=u_+-u_- =\sqrt{D^2-4fD},\qquad
 f=\frac{D^2-l^2}{4D}.
\]
这给出了用固定物屏间距与两次清晰成像位置测焦距的方法。$D=4f$时判别式为零，唯一位置$u=v=2f$、放大率$-1$；$D<4f$时没有满足实物实像要求的透镜位置。也可由$u+v\geq2\sqrt{uv}=2\sqrt{fD}$直接得到$D\geq4f$。实际厚透镜要以主面测距，若机械测的是两顶点位置，须先修正主面偏移。
\end{solution}

\begin{exercise}\label{ex:o1etendue}
理想近轴成像系统横向放大率为$m_y$，物像介质折射率分别$n_O,n_I$。求同一轴上物点光束的角放大率，并解释放大像时为何常需减小像方数值孔径。
\Diagram{fig:o1etendueq}{同一物点的两条光线定义孔径角}{\draw[axis](-.3,0)--(6,0);\draw[curve](2,-1.1)rectangle(3,1.1);\draw[curve](0,0)--(2,1);\draw[curve](0,0)--(2,-1);\draw[curve](3,.6)--(5,0);\draw[curve](3,-.6)--(5,0);}
\end{exercise}
\begin{solution}
图\ref{fig:o1etendueq}中，物像面之间使用约化变量$(y,p=n\theta)$。成像要求$B=0$，横向放大率为$A=m_y$。矩阵行列式恒为1，故$AD=1$、$D=1/m_y$，有
\[
 \begin{pmatrix}y_I\\n_I\theta_I\end{pmatrix}
 =\begin{pmatrix}m_y&0\\ C&1/m_y\end{pmatrix}
 \begin{pmatrix}y_O\\n_O\theta_O\end{pmatrix}.
\]
对轴上同一物点$y_O=0$，下行式化为$n_I\theta_I=n_O\theta_O/m_y$，所以角放大率为$n_O/(n_Im_y)$。
\Diagram{fig:o1etenduea}{像高和约化孔径角的乘积保持}{\draw[curve](0,0)rectangle(1.8,.7);\node at(.9,.35){$h_O$};\draw[curve](3,0)rectangle(5.7,1.1);\node at(4.35,.55){$h_I=m_yh_O$};\node[below]at(3,-.8){$n_Oh_O\theta_O=n_Ih_I\theta_I$};}
图\ref{fig:o1etenduea}所画两束使用同一组通光光线，其数值孔径满足$\mathrm{NA}_I\simeq\mathrm{NA}_O/|m_y|$。负放大率还带来倾角反号，孔径大小则取绝对值。对轴上附近小面元，面积放大$|m_y|^2$，小圆锥立体角约为$\pi\theta^2$，于是
\[
 n_I^2\dd A_I\dd\Omega_I
 =n_O^2\dd A_O\dd\Omega_O.
\]
该关系给出近轴通光相空间体积的不变性。无损时同一光束的功率也守恒，因此辐亮度除以$n^2$保持。扩大像面积会减小相应像方角锥；若另行增大物方收集孔径，所比较的已是另一组光线，不能把两种条件下的亮度混为一谈。
\end{solution}

\begin{exercise}\label{ex:o1magnification}
薄透镜在空气中$f>0$，物高$h$，物距$u>f$。求像距与横向放大率，并由两条光线画出$u=5f/3$的像。
\Diagram{fig:o1magnifyq}{已知物高、物距与焦距}{\OptAxis\draw[<->](2,-1.8)--(2,1.8);\draw[->,thick](0,0)--(0,1);\node[left]at(0,1){$h$};\fill(3.2,0)circle(1.3pt)node[below]{$F'$};\draw[<->](0,-.6)--(2,-.6)node[midway,below]{$u$};}
\end{exercise}
\begin{solution}
图\ref{fig:o1magnifyq}中$1/u+1/v=1/f$，解得$v=fu/(u-f)$。通过光心的近轴光线保持方向，由相似三角形$h'/v=-h/u$，所以横向放大率$m_y=h'/h=-v/u=-f/(u-f)$。
\Diagram{fig:o1magnifya}{平行轴光线经后焦点，过光心光线不偏折，两者交于倒立实像}{\OptAxis\draw[<->](2,-1.9)--(2,1.9);\draw[->,thick](0,0)--(0,1);\draw[->,thick](5,0)--(5,-1.5);\draw[->](0,1)--(2,1);\draw[->](2,1)--(5,-1.5);\draw[->](0,1)--(5,-1.5);\fill(3.2,0)circle(1.2pt)node[above]{$F'$};}
取$u=5f/3$，则$v=5f/2$、$m_y=-3/2$。图\ref{fig:o1magnifya}按$f=1.2$、$u=2$作图，实像在$z=5$、像高$-1.5h$，交点与计算吻合。真实透镜在大孔径或大物高时需考虑像差，本图只表达近轴模型。
\end{solution}

\begin{exercise}\label{ex:o1afocal}
正透镜$f_1>0$、$f_2>0$相隔$f_1+f_2$。求平行细光束通过后的束径比与小角偏转比，并说明无焦系统是否“没有作用”。
\Diagram{fig:o1afocalq}{共焦双正透镜系统}{\OptAxis\draw[<->](1,-1.4)--(1,1.4);\draw[<->](5,-1.1)--(5,1.1);\draw[guide](3.5,-1.6)--(3.5,1.6);\node[above]at(3.5,1.6){$F'_1=F_2$};}
\end{exercise}
\begin{solution}
图\ref{fig:o1afocalq}中$d=f_1+f_2$，直接乘矩阵得到
\[
 M=\begin{pmatrix}-f_2/f_1&f_1+f_2\\0&-f_1/f_2\end{pmatrix}.
\]
对于轴向平行光$\theta_i=0$，$y_o=-(f_2/f_1)y_i$，输出仍平行，束径绝对比$f_2/f_1$。
\Diagram{fig:o1afocala}{轴向光束先会聚再重新平行，图取两焦距不同}{\OptAxis\draw[<->](1,-1.4)--(1,1.4);\draw[<->](5,-1.4)--(5,1.4);\draw[->](0,1)--(1,1);\draw[->](1,1)--(5,-.6);\draw[->](5,-.6)--(6.5,-.6);\draw[->](0,-1)--(1,-1);\draw[->](1,-1)--(5,.6);\draw[->](5,.6)--(6.5,.6);}
对整体倾斜的小角平行光，$\theta_o=-(f_1/f_2)\theta_i$，与高度无关。图\ref{fig:o1afocala}取$f_1/f_2=5/3$，中间焦点位于第一透镜后2.5图示单位。无焦仅表示$C=0$、平行束仍为平行束；系统仍可扩束、缩束和改变角度，绝非恒等映射。
\end{solution}

\begin{exercise}\label{ex:o1glassplate}
一块厚度$d$、折射率$n$的平行玻璃板插入原本将于空气中会聚的近轴光束，板面垂直轴。求焦点沿轴的位移，忽略反射。
\Diagram{fig:o1glassq}{会聚光束在平行板内传播的斜率减小}{\OptAxis\draw[curve](2,-1.1)rectangle(3,1.1);\draw[guide](0,1.3)--(4.8,0);\draw[guide](0,-1.3)--(4.8,0);\node[above]at(2.5,1.1){$n,d$};\fill(4.8,0)circle(1.3pt)node[below]{$F_0$};}
\end{exercise}
\begin{solution}
图\ref{fig:o1glassq}中原光线以斜率$\theta<0$会聚。进入平面板后$\theta_g=\theta/n$，传播$d$导致高度变化$d\theta/n$；若仍为空气则变化$d\theta$。出板后倾角恢复$\theta$，故高度比原光线高出$\Delta y=d\theta(1/n-1)>0$。
\Diagram{fig:o1glassa}{出板后两光线平行，新的交轴点向后移动}{\OptAxis\draw[curve](2,-1.1)rectangle(3,1.1);\draw[->](0,1.2)--(2,.7)--(3,.575)--(5.3,0);\draw[guide](2,.7)--(4.8,0);\fill(4.8,0)circle(1.2pt)node[below]{$F_0$};\fill(5.3,0)circle(1.2pt)node[above right]{$F$};}
图\ref{fig:o1glassa}中的轴向位移由$\theta\Delta z=-\Delta y$得$\Delta z=d(1-1/n)>0$。矩阵角度看，是把空气传播长度$d$的约化距离$d$替换为$d/n$，而折射后的外部倾角不变。大孔径时不同光线位移并不完全相同，会引入球差。
\end{solution}

\begin{exercise}\label{ex:o1turning}
水平分层介质$n(y)=n_0\sqrt{1-y^2/a^2}$。光线从$y=0$以与$+z$轴成$\alpha_0$的方向进入。用守恒量求最大高度，不使用近轴条件。
\Diagram{fig:o1turnq}{光线在折射率下降区弯折至水平}{\OptAxis\draw[axis](0,0)--(0,2.7)node[above]{$y$};\draw[curve](0,0)..controls(1,1)and(2,1.8)..(3.2,1.8);\draw[guide](0,1.8)--(4,1.8);\node[left]at(0,1.8){$y_t$};\node at(.8,.3){$\alpha_0$};}
\end{exercise}
\begin{solution}
图\ref{fig:o1turnq}中，介质沿$z$平移不变，$b=n\cos\alpha=n_0\cos\alpha_0$守恒。用$y'=\tan\alpha$表示方向，得到
\[
 n^2(y)=b^2(1+y'^2),\qquad
 y'^2=\frac{\sin^2\alpha_0-y^2/a^2}{\cos^2\alpha_0}.
\]
真实射线要求右侧非负，允许区域为$|y|\leq a|\sin\alpha_0|$。最高点$y'=0$，所以正向上射线的最大高度为$y_t=a\sin\alpha_0$；若只讨论幅度，则为$a|\sin\alpha_0|$。
\Diagram{fig:o1turna}{允许高度由$n(y)\geq n_0\cos\alpha_0$确定}{\draw[axis](0,0)--(4.2,0)node[right]{$y/a$};\draw[axis](0,0)--(0,2.8)node[above]{$n/n_0$};\draw[curve,domain=0:.99,samples=70]plot({3.5*\x},{2.3*sqrt(1-\x*\x)});\draw[guide](0,1.7)--(3.7,1.7);\node[left]at(0,1.7){$\cos\alpha_0$};}
图\ref{fig:o1turna}中的折射率曲线与水平$b$线相交处就是转点。还可将方程积分，以核对近轴极限。对$0<\alpha_0<\pi/2$的上升支，令$A=a\sin\alpha_0$，则
\[
 \int_0^{y}\frac{\dd q}{\sqrt{A^2-q^2}}
 =\frac1{a\cos\alpha_0}\int_0^z\dd s,\qquad
 [\arcsin(q/A)]_0^{y}=\frac{z}{a\cos\alpha_0}.
\]
这里用$q=A\sin u$，在上升支$0\leq u\leq\pi/2$有$\dd q=A\cos u\dd u$、$\sqrt{A^2-q^2}=A\cos u$。反解得$y=A\sin[z/(a\cos\alpha_0)]$，在转点处接续下降支。
精确轴向周期是$2\pi a\cos\alpha_0$，最高点在$z_t=\pi a\cos\alpha_0/2$。小角时振幅$A\simeq a\alpha_0$、周期趋$2\pi a$，恢复近轴解。有限角度下转点高度与周期都应使用轴向守恒量求出，不能仅保留振幅修正而仍把近轴周期称为精确周期。
\end{solution}

\begin{exercise}\label{ex:o1centralturn}
圆对称介质在某区域$n(r)=n_a a/r$。求一条非径向光线能否有孤立最近点，推导其轨迹形状。讨论守恒量$b=n_a a$的特殊情形。
\Diagram{fig:o1centralq}{径向折射率场中的光线与径向夹角}{\draw[curve](0,0)circle(1);\draw[curve](0,0)circle(2);\draw[axis](0,0)--(2,1)node[right]{$r$};\draw[axis](2,1)--(2.5,2.2)node[above]{切向};\node[below left]at(0,0){$O$};}
\end{exercise}
\begin{solution}
图\ref{fig:o1centralq}的$n(r)r=n_aa$为常数，所以$\sin\psi=b/(n_aa)$也为常数。若$0<b<n_aa$，夹角固定且不等于直角，不存在孤立转点。由$\dd r/\dd s=\cos\psi$、$r\dd\varphi/\dd s=\sin\psi$，得$\dd\varphi/\dd r=\tan\psi/r$。
\Diagram{fig:o1centrala}{固定夹角对应对数螺线，图只画有限介质区域}{\draw[axis](-2.3,0)--(2.5,0);\draw[axis](0,-2.3)--(0,2.5);\draw[curve,domain=0:300,samples=100]plot({.35*exp(.006*\x)*cos(\x)},{.35*exp(.006*\x)*sin(\x)});\node[below left]at(0,0){$O$};}
以$(r_0,\varphi_0)$为初始点积分，$\int_{\varphi_0}^{\varphi}\dd\phi=\tan\psi\int_{r_0}^{r}\dd q/q=\tan\psi[\ln q]_{r_0}^{r}$，得$\varphi-\varphi_0=\tan\psi\ln(r/r_0)$，即图\ref{fig:o1centrala}的对数螺线。向内支的$\cos\psi<0$，须保留$\tan\psi$的符号。若$b=n_aa$，则$\psi=\pi/2$，光线可保持任意给定半径的圆周轨道；此时不能再除以$\dd r/\dd s$，应直接使用切向关系。该折射率在$r=0$奇异，轨迹仅适用于题定介质区域。
\end{solution}

\begin{exercise}\label{ex:o1phasors}
同频同偏振两束光复振幅为$A$和$2A\ee^{i\delta}$。求合振幅的大小、相位以及最大最小强度，并说明相位求值应怎样处理象限。
\Diagram{fig:o1phasorq}{两相量大小之比为1:2}{\draw[axis](0,0)--(1.5,0)node[right]{$A$};\draw[axis](0,0)--(1.5,2.598076)node[right]{$2A\ee^{i\delta}$};\draw(.6,0)arc(0:60:.6);\node at(.8,.4){$\delta$};}
\end{exercise}
\begin{solution}
以第一束的相位为零，图\ref{fig:o1phasorq}中合复振幅为
\[
 \frac{\mathcal E}A=1+2\ee^{i\delta}
 =1+2\cos\delta+2i\sin\delta.
\]
因此$|\mathcal E|^2/A^2=(1+2\cos\delta)^2+4\sin^2\delta=5+4\cos\delta$，合振幅大小为$A\sqrt{5+4\cos\delta}$，相位为
\[
 \phi=\operatorname{atan2}(2\sin\delta,1+2\cos\delta).
\]
该函数由实、虚两分量共同确定象限；单独取反正切商只能给出相差整数$\pi$的一组可能值。
\Diagram{fig:o1phasora}{合相量实部可能为负，单值反正切会丢失象限}{\begin{scope}[scale=.7]\draw[axis](-2.2,0)--(3.4,0)node[right]{实部};\draw[axis](0,-2.3)--(0,2.8)node[above]{虚部};\draw[curve](1,0)circle(2);\draw[->,thick](0,0)--(-.5,1.322876);\node[left]at(-.5,1.322876){$\mathcal E$};\end{scope}}
图\ref{fig:o1phasora}的合相量末端以$(A,0)$为圆心、$2A$为半径。原点位于圆内，合相量会经过左半平面，但不会为零，因为两束振幅不等。若$A$对应单束强度$I_1$，则
\[
 I_{\max}=9I_1\quad(\delta=2m\pi),\qquad
 I_{\min}=I_1\quad(\delta=(2m+1)\pi),
\]
可见度为$(9-1)/(9+1)=4/5$。合相量越过虚轴、实部变号时，相位连续经过$\pm\pi/2$；越过负实轴时，取值区间为$(-\pi,\pi]$的主值相位才会发生$2\pi$跳变。本题在$\delta=\pi$通过负实轴，可按连续相位消去这种记号上的跳变。场本身始终连续且不为零，相位也始终有定义。
\end{solution}
